\textbf{5.\ \textsc{Force sensor} (5 points)} --- \textit{Päivo Simson.}

A force sensor is constructed from a flexible beam of working length $L$ and
height $h$ to which four identical electrically resistive wires of length
$l \ll L$ are attached, as shown in the figure below. If a force $F$ is applied
to the end of the beam, it will bend the beam and stretch the upper wires and
compress the lower ones, causing changes in the electrical resistance. Using
the Wheatstone bridge method, these changes can be converted into a voltmeter
reading, allowing for the measurement of the applied force $F$.

\begin{center}\includegraphics[width=0.41\linewidth]{figures/NBPhO_2023_Q5_fig1.png}\end{center}

In what follows, assume that the deflection of the beam is very small.

\textbf{i)} \textit{(2 points)} The bending moment $M(x)$ (i.e. the torque with
which one fictitious half of the beam affects the other) is inversely
proportional to the curvature radius $r(x)$ of the center line of the beam:
$M = EI/r$, where $E$ is Young's modulus and $I$ is a constant depending on the
geometry of the of the beam's cross-section (both are known). Find the
elongation $\Delta l$ of the upper wire when a force $F$ is applied to the
sensor. Assume that the wires are placed in the middle of the beam.

\textbf{ii)} \textit{(1 point)} Find the resistances $R_1$ and $R_2$ when
$\Delta l$ from the previous section is known and the initial resistance of the
wires is $R_0$. Assume that the volume of the wire remains constant during the
deformation.

\textbf{iii)} \textit{(2 points)} The resistors are arranged in a Wheatstone
bridge configuration shown in the figure below, where $U$ is the known battery
voltage. Find the relationship between the measured voltage $V$ and the force
$F$.

\begin{center}\includegraphics[width=0.42\linewidth]{figures/NBPhO_2023_Q5_fig2.png}\end{center}
